\documentclass[11pt]{article}

\usepackage[margin=2.5cm]{geometry}
\usepackage[utf8]{inputenc}
\usepackage[T1]{fontenc}
\usepackage{enumitem}
\usepackage{hyperref}
\usepackage{titling}
\usepackage{amsmath}
\usepackage{booktabs}

\setlist{itemsep=2pt, topsep=4pt}

\title{Enigma Engine \\ \large Project Presentation and User Manual}
\author{Interactive Graphics -- Final Project}
\date{}

\begin{document}

\maketitle
\tableofcontents

\section{Introduction}

Enigma Engine is a browser-based, real-time 3D car experience built with
WebGL. It lets the user drive a car around a race track using a
physics-based driving model, customize the car's paint and materials, open
its doors and windows, control its lights, and switch between several
camera views (orbit, free camera, driver view, drone view). A dashboard
sidebar exposes all of these options through a simple graphical interface,
and a heads-up display shows live telemetry (speed, RPM, gear, lap time).

Everything runs directly in the browser: there is no installation step, no
backend server, and no build/compile step. This document is both a
technical presentation of how the project is built and a short user manual
explaining how to use it.

\section{Environment}

\subsection{Runtime and module loading}
The project is built on top of \textbf{WebGL}, accessed through the
\textbf{Three.js} rendering library (version 0.160.0). The application is a
single HTML page (\texttt{index.html}) that loads plain JavaScript ES
modules directly, with no bundler (Webpack, Vite, Parcel, ...) and no
package manager (no \texttt{node\_modules}, no \texttt{npm install}). This
keeps the project extremely easy to run: opening \texttt{index.html} in a
local static server is enough.

All third-party libraries are declared once, in the page's \emph{import
map}, and are downloaded from public content delivery networks (unpkg and
jsDelivr) the first time the page is opened:

\begin{verbatim}
"three":            "https://unpkg.com/three@0.160.0/build/three.module.js"
"three/addons/":    "https://unpkg.com/three@0.160.0/examples/jsm/"
"@tweenjs/tween.js":"https://unpkg.com/@tweenjs/tween.js@23.1.1/..."
"cannon-es":        "https://cdn.jsdelivr.net/npm/cannon-es@0.20.0/+esm"
\end{verbatim}

Every JavaScript file in the project then simply writes
\texttt{import * as THREE from 'three'} (or similar), exactly as it would
with a bundler, but without any local installation step.

\subsection{Rendering pipeline configuration}
The renderer is configured for a realistic look:
\begin{itemize}
    \item Antialiasing enabled, \texttt{powerPreference: "high-performance"}
    to request the discrete GPU when available.
    \item Output color space set to sRGB, with ACES filmic tone mapping
    (exposure 1.0) for a more cinematic contrast curve.
    \item Shadows enabled with a soft PCF shadow map.
    \item Pixel ratio capped at 1.5 (instead of the raw device pixel ratio)
    to keep the frame rate stable on high-DPI screens.
\end{itemize}
The camera is a standard perspective camera with a 60\textdegree{} field of
view and near/far clipping planes at 0.1 and 1000 units.

\subsection{Application startup sequence}
When the page loads, the application: (1) creates the scene, camera and
renderer; (2) loads the track and the car model in parallel; (3) builds the
physics world and the car's rigid body from the loaded meshes; (4) wires up
every UI control (lights, colors, animations, camera buttons, HUD); (5)
runs a short \emph{prewarm} pass that forces all shaders to compile and
renders a few warm-up frames (both in day and night lighting) so that the
GPU pipeline is fully "hot" before anything is shown; and finally (6) hides
the loading overlay and starts the real render loop. This prewarm step
avoids a visible stutter on the very first frames, which is otherwise a
common issue with WebGL applications that compile shaders lazily.

\section{Libraries, Tools and Models Not Developed by the Team}

The following third-party software and assets are used in the project but
were not created by the team:

\begin{itemize}
    \item \textbf{Three.js} (v0.160.0) -- WebGL rendering engine, together
    with its official add-ons: \texttt{GLTFLoader} and \texttt{DRACOLoader}
    (to load the compressed 3D models) and \texttt{OrbitControls} (to orbit
    the camera around the car).
    \item \textbf{Draco} -- Google's mesh compression decoder (hosted on
    Google's own CDN), used by \texttt{GLTFLoader} to decompress the 3D
    model geometry when it is loaded.
    \item \textbf{tween.js} (v23.1.1) -- animation/tweening library used to
    smoothly animate doors, windows, lights intensity and camera
    transitions between two states.
    \item \textbf{cannon-es} (v0.20.0) -- physics engine used for the whole
    vehicle simulation: rigid bodies, gravity, suspension, wheel raycasting
    and collisions with the track.
    \item \textbf{Car model} (\texttt{car.glb}, \textasciitilde18\,MB) -- a
    detailed car model obtained from an external source, not modeled by the
    team; its exported material file still references the original
    Blender project and artist folder names.
    \item \textbf{Track model} (\texttt{track\_1.glb}, \textasciitilde69\,MB)
    -- the race circuit and its surrounding environment (road, trees, lamp
    posts, garage), also an external asset.
    \item \textbf{Textures} -- the day/night equirectangular sky panoramas
    and the PBR ground/floor textures (diffuse, normal and roughness maps)
    are free texture packs, not created by the team.
    \item \textbf{Audio samples} -- the engine sound recordings (five
    samples captured at fixed RPM values, from idle to redline) and the
    short sound effects (door open/close, turn signal click, engine
    startup) are external audio files, not recorded by the team.
\end{itemize}

Everything else -- the whole application logic, the physics setup, the
lighting and reflection system, the UI, the audio synthesis, and the way
all of the above assets are combined -- was written by the team.

\section{Project Structure}

The repository is organized into two top-level folders:

\begin{verbatim}
final-project-enigma/
|-- code/
|   |-- page/
|   |   |-- index.html      # entry point, import map, UI markup
|   |   \-- style.css       # UI/HUD styling
|   \-- script/             # ES modules, no bundler
|       |-- animations.js   # part animation + click-to-interact system
|       |-- audio.js        # Web Audio SFX and engine sound synthesis
|       |-- bestLap.js      # lap timing via a trigger zone
|       |-- camera.js       # orbit / free / driver / drone camera modes
|       |-- car_model.js    # declarative description of animatable parts
|       |-- color.js        # material registry, live color/roughness edit
|       |-- engine.js       # engine and gearbox simulation
|       |-- environment.js  # track loading, sky textures, tree instancing
|       |-- lights.js       # scene lights, car lights, street lamp pool
|       |-- loaders.js      # shared GLTFLoader + DRACOLoader
|       |-- main.js         # bootstraps the app, render/animate loop
|       |-- model.js        # generic GLTF loader + per-part tweening
|       |-- physics.js      # vehicle physics (RaycastVehicle)
|       |-- reflections.js  # dynamic cube-map reflections
|       |-- scene.js        # scene, camera and renderer setup    
|       |-- settings.js     # small global settings store
|       \-- ui.js           # DOM event wiring for sidebar and HUD
\-- src/
    |-- models/             # car and track GLB files
    |-- textures/           # sky and PBR ground/floor textures
    \-- audio/               # engine RPM samples and short SFX
\end{verbatim}

Each JavaScript module has a single responsibility, and \texttt{main.js}
only wires the modules together: it does not itself contain rendering,
physics or UI logic.

\section{Technical Aspects}

\subsection{Scene, camera and the render loop}
The scene uses a perspective camera and the WebGL renderer described in
Section~2.2. The main render loop runs on
\texttt{requestAnimationFrame}. On every frame it, in order:

\begin{enumerate}
    \item advances the physics simulation by one step;
    \item updates the on-screen speed / RPM / gear readouts;
    \item updates the procedural engine sound based on the new RPM;
    \item moves the sun light so it stays above and around the car;
    \item updates the active camera mode (orbit, free, driver or drone);
    \item advances any running tween animation (doors, lights, transitions);
    \item refreshes the car's dynamic reflections and the nearby street
    lamps;
    \item checks the lap-timer trigger zone against the car's position;
    \item renders the frame.
\end{enumerate}

A small FPS counter (updated once per second) is shown in the corner for
debugging and performance monitoring.

\subsection{Vehicle physics}
The car is simulated as a rigid body using \texttt{cannon-es}'s
\texttt{RaycastVehicle}: a box-shaped chassis with four independently
suspended wheels. Table~\ref{tab:wheels} summarizes the main tuning
parameters.

\begin{table}[h]
\centering
\caption{Wheel and suspension parameters}
\label{tab:wheels}
\begin{tabular}{lcc}
\toprule
Parameter & Front wheels & Rear wheels \\
\midrule
Radius (m) & 0.338 & 0.345 \\
Suspension rest length (m) & \multicolumn{2}{c}{0.22} \\
Suspension stiffness & 35 & 38 \\
Damping (compression / relaxation) & 2.5 / 3.8 & 2.3 / 3.6 \\
Friction slip & 4.5 & 4.5 \\
Max steering angle & 25\textdegree{} & -- \\
\bottomrule
\end{tabular}
\end{table}

The chassis is a box of roughly $0.90 \times 0.60 \times 2.25$ metres with a
mass of 1500\,kg, spawned at a fixed point on the track. Gravity is set to
$-9.82\,\mathrm{m/s^2}$ and the physics world steps at a fixed 60\,Hz with
sub-stepping, independently of the rendering frame rate.

The track geometry is converted into a static physical collision mesh: at
load time, every track mesh that is not a tree is turned into an exact
\texttt{Trimesh} collider, so the car collides with the real shape of the
road, curbs and walls instead of a simplified approximation.

Several extra behaviors sit on top of the base physics simulation:
\begin{itemize}
    \item \textbf{Smoothed steering}: the steering angle does not snap
    instantly to the key input, it eases toward its target and back to
    center, capped at $\pm$25\textdegree{}.
    \item \textbf{Speed-dependent understeer}: above roughly 50\,km/h, if
    the requested steering angle is large relative to the current speed,
    the front wheels' grip and effective steering angle are both reduced,
    so the car cannot corner unrealistically sharply at high speed.
    \item \textbf{Low-speed brake lock}: when the car is nearly stopped and
    the brake is held with no throttle, its remaining velocity is zeroed
    directly, so it does not creep forever at a tiny residual speed.
    \item \textbf{Rollover recovery}: if the car's "up" direction tips more
    than about 101\textdegree{} from vertical (i.e.\ it has flipped), it
    is automatically lifted, its roll and pitch are reset to level (yaw is
    kept), and its velocity is cleared, so the player is never permanently
    stuck upside down.
\end{itemize}

\subsection{Engine and transmission model}
On top of the physics engine, a hand-written engine model computes the
wheel force from a torque curve and a gearbox, instead of relying on a
single generic acceleration value. The torque curve is a lookup table,
linearly interpolated between the sampled points shown in
Table~\ref{tab:torque}.

\begin{table}[h]
\centering
\caption{Torque curve (RPM $\rightarrow$ Nm)}
\label{tab:torque}
\begin{tabular}{ccccccc}
\toprule
RPM & 800 & 1500 & 2000 & 3000 & 4000 & 5000 \\
Nm & 220 & 260 & 290 & 350 & 400 & 440 \\
\midrule
RPM & 6000 & 7000 & 8000 & 8500 & 9000 & 9200 \\
Nm & 470 & 460 & 430 & 390 & 340 & 150 \\
\bottomrule
\end{tabular}
\end{table}

The gearbox has six forward gears plus reverse, with the ratios in
Table~\ref{tab:gears} (all multiplied by a fixed final-drive ratio of
3.96):

\begin{table}[h]
\centering
\caption{Gear ratios}
\label{tab:gears}
\begin{tabular}{lcccccc}
\toprule
Gear & 1 & 2 & 3 & 4 & 5 & 6 \\
\midrule
Ratio & 4.75 & 3.08 & 2.40 & 2.04 & 1.68 & 1.18 \\
\bottomrule
\end{tabular}
\end{table}

The gearbox is fully automatic: it shifts up when the engine goes above
about 7000\,RPM, and shifts down when RPM drops below a threshold that also
depends on how much throttle is applied (so a light touch of the throttle
does not trigger an unwanted downshift). A short cooldown between shifts
prevents rapid gear "hunting". Idle RPM is 800 and redline is 9000; engine
RPM itself is smoothed with different response speeds when rising versus
falling, to mimic engine inertia. The resulting wheel force is:
\[ F_{\text{wheel}} = \frac{\text{torque}(RPM) \times \text{gear ratio}
\times \text{final drive} \times \text{efficiency}}{\text{wheel radius}} \]
which is applied directly to the rear wheels of the \texttt{RaycastVehicle}.

\subsection{Lighting system}
The scene has a directional "sun" light (intensity 8.5, casting a
$2048\times2048$ shadow map) which is re-targeted every frame to stay above
the car, keeping shadow quality high near the player regardless of the
track size. An ambient hemisphere light (intensity 1.2) provides soft
indirect sky/ground lighting.

A day/night cycle smoothly interpolates a 0--24h virtual clock over about
two seconds when the "Night Mode" switch is toggled. This clock drives: the
sky texture (day or night equirectangular panorama), the sun and hemisphere
light intensities (eased with a smoothstep curve), and whether the street
lamps are active.

Each of the car's own lights (low/high beams, running lights, tail lights,
brake lights, turn signals, interior ambient light) is built by locating
the corresponding named mesh inside the 3D model, upgrading its material so
it can glow (\texttt{emissive} property), and attaching a matching spot
light. Turning a light on or off is animated with a short tween instead of
an instant on/off, for a more natural look. Turn signals additionally blink
on a fixed interval and play a click sound on every "on" phase.

Street lamps deserve a special mention: the track model can contain dozens
of lamp posts, but creating one real light per lamp post would be too
expensive. Instead, a fixed pool of 21 spotlights is created once; every
frame, the lamp posts are sorted by distance to the car and the pool is
re-assigned to the 21 closest ones. This keeps the lighting cost constant
regardless of how large the track is.

\subsection{Reflections}
The car's paint reacts to its surroundings using a dynamic cube-map
reflection: a \texttt{CubeCamera} renders the environment from a point just
above the car into a small (64px) cube render target, which is then used as
the environment map for the car's materials. Rendering a full cube map
every frame would be expensive, so it is refreshed only once every three
frames by default, and a "Static Reflections" setting lets the user freeze
it completely for better performance on slower machines.

\subsection{Model, materials and part-animation system}
The car's animatable parts (doors, windows, hood, rear wing, ignition key,
wheels, steering wheel) are described declaratively in a single
configuration object, rather than being hard-coded one by one: each part
lists which mesh it corresponds to, how it should rotate or move, how long
the animation should take, which sidebar switch controls it, and which
sound effects to play. A generic loader reads this configuration once the
3D model is loaded and precomputes the "closed" and "open" transform for
every part, so that later toggling only needs to interpolate between two
already-known states.

Materials are collected into a name-based registry as the car model is
loaded (using the material names authored in Blender), which is what
allows the sidebar's color pickers and metallic/roughness sliders to edit
the live materials directly.

Clicking on a part in the 3D scene (instead of using the sidebar) works
through raycasting: only the meshes that belong to a clickable part are
considered, both for detecting hover (to draw a highlight) and for
detecting clicks, which keeps this check fast even though the whole car
model has many more meshes than just the clickable ones.

\subsection{Audio system}
All audio uses the native Web Audio API directly, with no external audio
library. Short one-shot effects (door open/close, turn signal click, engine
startup) are simply decoded once and played back on demand.

The running engine sound is synthesized rather than looped from a single
clip: five samples were recorded at fixed RPM values (900, 2294, 4139,
6025, 8198), and at runtime the two samples that bracket the current RPM
are cross-faded using an equal-power (cosine-based) curve, while each
sample's playback rate is also adjusted so its pitch matches the exact
current RPM. A second pair of gain buses blends between an "on-throttle"
and an "off-throttle / coasting" mix depending on how much the accelerator
is pressed, so the engine also sounds different when accelerating versus
coasting at the same RPM.

\subsection{Lap timing}
Lap timing uses an invisible trigger volume placed at a "finish line"
marker inside the track model. Every frame, the car's world position is
converted into the trigger's local space and compared against its bounds;
crossing into the zone for the first time starts the lap clock, and
crossing it again on a later lap stops the current lap and compares it
against the stored best time, updating the on-screen best lap
(\texttt{mm:ss:ms}) if the new lap was faster.

\subsection{Performance optimizations -- summary}
Several techniques throughout the project exist purely to keep the frame
rate high:
\begin{itemize}
    \item \textbf{GPU instancing} for trees: every tree of the same type on
    the track is drawn with a single \texttt{InstancedMesh} draw call
    instead of one draw call per tree, with small per-instance random
    rotation/scale for visual variety.
    \item \textbf{Light pooling} for street lamps (Section~5.4): a fixed
    number of real lights is reused for however many lamp posts exist.
    \item \textbf{Throttled and freezable reflections} (Section~5.5): the
    cube-map reflection is updated every few frames, or not at all in
    static mode.
    \item \textbf{Filtered raycasting} (Section~5.6): click/hover detection
    only tests the small set of clickable meshes, not the whole car model.
    \item \textbf{Shader/pipeline prewarming} (Section~2.3): shader
    compilation and the first heavy frames happen behind the loading
    screen, not during actual gameplay.
\end{itemize}

\section{User Manual -- Implemented Interactions}

\subsection{Getting started}
Open \texttt{code/page/index.html} through a local web server (opening the
file directly may be blocked by the browser's security policy for module
imports). Wait for the loading screen to disappear, then press the
\textbf{ENGINE} button in the instrument cluster to start the car, select
\textbf{D} in the gear selector, and use the arrow keys to drive.

\subsection{Driving controls}
\begin{table}[h]
\centering
\caption{Driving controls}
\begin{tabular}{ll}
\toprule
Input & Action \\
\midrule
Up arrow & Accelerate \\
Down arrow / Space & Brake \\
Left / Right arrow & Steer \\
Down arrow / S & Brake lights on (while held) \\
HUD engine button & Start / stop the engine (with crank sound) \\
HUD gear selector (N / D / R) & Neutral / Drive (automatic) / Reverse \\
\bottomrule
\end{tabular}
\end{table}
The engine only actually starts once the crank sound has finished playing,
so there is a short delay between pressing the button and being able to
move.

\subsection{Camera controls}
\begin{table}[h]
\centering
\caption{Camera controls}
\begin{tabular}{ll}
\toprule
Input / button & Action \\
\midrule
Mouse drag (orbit mode) & Rotate the camera around the car \\
Mouse scroll (orbit mode) & Zoom in / out (clamped range) \\
W / A / S / D, Q / E (free mode) & Fly forward/back/left/right, down/up \\
Mouse drag (free / driver mode) & Look around \\
"Orbit Camera" / "Free Camera" button & Toggle between orbit and free-fly \\
Compass buttons (Front/Back/Left/Right/Top) & Jump to a preset view \\
"Driver View" button & First-person onboard camera \\
"Drone View" button & Top-down camera following the car \\
\bottomrule
\end{tabular}
\end{table}
Switching to a preset compass view plays a smooth, eased transition instead
of an instant camera cut. Driver view lets the player look around with the
mouse; releasing the mouse smoothly re-centers the view.

\subsection{Car customization panel}
The sidebar "Dashboard" is organized into collapsible sections:
\begin{itemize}
    \item \textbf{Color Customization}: a color picker plus metallic and
    roughness sliders for the body paint, rims, brake calipers, seat fabric
    and steering wheel, applied live to the 3D model.
    \item \textbf{Animations}: switches to open/close the hood, rear wing,
    and both doors and windows. The same parts can also be opened by
    clicking on them directly in the 3D scene; hovering over a clickable
    part highlights it in blue.
    \item \textbf{Lights}: switches for running lights, low beams, high
    beams, left/right turn indicators, hazard lights and the interior
    ambient light. These switches follow realistic real-car dependencies:
    turning on the high beams automatically turns on the low beams, and
    turning off the running lights also turns off both beam types.
    \item \textbf{Settings}: a switch to enable/disable the blue hover
    highlight on clickable parts, and a "Static Reflections" switch that
    freezes the car's cube-map reflections to save performance.
    \item \textbf{Time of Day}: a single "Night Mode" switch that smoothly
    transitions the whole scene between day and night lighting over about
    two seconds.
\end{itemize}

\subsection{Heads-up display}
The main view overlays a small instrument cluster and telemetry panel:
\begin{itemize}
    \item \textbf{Speed} in km/h.
    \item \textbf{RPM}, shown both as a number and as a color-coded rev
    bar (green $\rightarrow$ orange $\rightarrow$ red, flashing near the
    redline).
    \item \textbf{Gear} readout (current numeric gear in Drive, or N / R).
    \item \textbf{Best lap} time, in \texttt{mm:ss:ms} format, updated
    automatically whenever a faster lap is completed.
    \item \textbf{FPS counter}, updated once per second.
\end{itemize}

\section{Known Limitations}
A few assets bundled in the repository are not currently used by the
application (an unused garage/box model and a couple of unreferenced
ground textures), left over from earlier iterations of the project; they
have no effect on the running application. There is no menu, pause screen
or multiplayer mode -- the experience is a persistent, single-player
free-driving and lap-timing session, which matches the scope of the
project.

\section{Conclusion}
Enigma Engine combines a real-time WebGL renderer, a raycast-based vehicle
physics simulation with a custom engine/gearbox model, a dynamic day/night
lighting system, procedurally synthesized engine audio, and a fully
interactive customization UI, all running client-side in the browser with
no build step. The result is an interactive car showroom and driving
simulator that can be opened and used directly from a static web page.

\end{document}
